\section{Economic Motivation}

Risky-asset prices reflect both expected cash flows and the discount rates
applied to those cash flows. Lettau and Ludvigson construct \cay{} as a
theory-motivated macroeconomic state variable intended to reveal variation
in expected returns and the equity risk premium \citep{lettau_ludvigson_2001}.

\subsection{Present-Value Relation}

The household intertemporal budget constraint implies that variation in the
consumption-to-total-wealth ratio must be associated with predictable future
returns, predictable future consumption growth, or both. The empirical
challenge is that total wealth includes unobservable human capital.

\subsection{Risk-Premium Interpretation}

The paper uses labor income to capture the long-run component of human
capital and estimates a stationary trend deviation among consumption, asset
wealth, and labor income. A high value of \cay{} is associated with high
expected excess returns in the original sample. This is consistent with a
countercyclical risk premium, but \cay{} is not a direct estimate of risk
aversion.
